\section{Cross-Family Results and Grounding}
\label{sec:crossfamily}

\begin{center}
\small
\begin{tabular}{lrrr}
\toprule
Input & R & A & B \\
\midrule
N2 & 69/75 & 62/75 & 69/75 \\
CS1 & 19/20 & 16/20 & 16/20 \\
\bottomrule
\end{tabular}

\captionof{table}{Qwen \texttt{qwen3.8-max-0902} on allocated
determinate opportunities. Of 42 unique first submissions, 39 returned
complete responses and 38 satisfy the output contract. Transport and
Schema failures remain in these denominators.}
\label{tab:qwen}
\end{center}

\subsection{A model--interface comparison}
Qwen receives the same 14 buildings and R/A/B core tasks through a
tool-free multi-image API. Thinking is enabled with a 32,768-token
budget, completion is capped at 49,152 tokens, and temperature is 0.6.
Its R wrapper additionally contains the generic world Schema, whereas
GPT-R does not. This Schema contains no building answers.
GPT can crop assigned local pixels; Qwen uses submitted image inputs.
The settings and tools are not compute-matched.

Two initial non-streaming requests disconnected without final output.
They remain failures, not retried replacements.
The remaining 40 tasks each received one first streaming submission;
one also disconnected. The comparison includes the actual delivery
conditions, not just successful replies.

\subsection{N2 mostly aligns answers and worlds}
In N2, R is correct on 69/69 contract-valid opportunities, A on 62/63,
and B on 69/70. Of 63 common-valid R/A opportunities, 62 are jointly
correct and one is R-only. A/B share 58 valid opportunities:
57 jointly correct and one jointly wrong, with no rescue or regression.
The seven-readout increase on the full denominator from A to B comes
from different failure sets, not paired content improvement.

The single R/A mismatch is a school stair connection whose generated
endpoint interface is a corridor, not a landing. This is a contract-level
instantiation failure, not established evidence of wrong physical
identity. The strong N2 alignment shows that the observed
answer--construction difference varies across model--interface conditions.

\subsection{CS1 exposes stair implementation failures}
All six CS1 outputs are contract-valid. R/A/B yield 19/20, 16/20,
and 16/20. All eight determinate identity relations per route are
correct. Three A readouts fail despite correct R answers:
a Lovell landing declares \texttt{floor\_z}=5.29\,ft while its
material top is 5.79\,ft; a Gamble assembly's material differs from
the union of its parts; another endpoint directly uses a corridor
instead of a landing.

B repairs the two-riser floor relation that R and A both missed,
but damages a previously correct endpoint relation. A/B thus have
one rescue and one regression despite equal totals.
No definite material obstruction is detected in the 25 valid Qwen
construction worlds.

\subsection{Grounding and insufficient evidence}
Grounding links world instances back to the assigned drawings.
For GPT N2, given-point box IoU$\geq0.25$ holds in 446/656 A records
and 440/656 B records (mean IoU 0.491/0.489); 16 abstract-region
references per route are not evaluable. CS1 yields 49/88 and 40/88.
This is not full-page detection or complete 3D grounding.

Strict source-ID matching makes some Qwen boxes score zero:
56 existing boxes use aliases such as \texttt{S\_V03} for
\texttt{V03}. Under a lexical-alias diagnostic assumption, 55 reach
the IoU threshold. This does not replace primary scores.
Other failures lack target boxes or have low overlap; low IoU alone
cannot establish that an object was not seen.

On unresolved identity, GPT N2 preserves \unk{} only 2/24 times,
versus 18/24 in v1. Qwen R/A assert identity and B fails to return:
two unsupported certainties plus a service failure, not three proven
identity errors.
